\section{Protocolo experimental}

\subsection{Condiciones comunes}

Todos los experimentos usan los mismos 12,753 ejemplos de entrenamiento y 3,712
de validación de las 104 clases, la semilla 0 y lotes de 64 imágenes. La
\emph{referencia} es la CNN de cuatro convoluciones elegida en el Avance~4:
dropout de 0.5 antes de la salida, aumento de datos solo en entrenamiento
(volteo horizontal, rotación, zoom y contraste), AdamW con decaimiento de pesos
de 0.01 \autocite{loshchilov2019adamw}, tasa inicial de $10^{-3}$ reducida a la
mitad cuando la pérdida de validación no mejora en dos épocas (mínimo
$10^{-5}$), entrada de $112\times112$ y hasta 150 épocas. El entrenamiento se
detiene tras 12 evaluaciones sin mejorar el F1 macro de validación y se conserva
el checkpoint con mayor F1 macro; todas las métricas reportadas corresponden a
ese checkpoint, no a la última época.

\subsection{Un cambio a la vez}

Para atribuir cada diferencia a una causa, cada experimento modifica un solo
factor respecto a la referencia: la arquitectura (MLP o CNN ampliada), la pérdida (pesos por clase), la regularización o un
hiperparámetro (tasa de aprendizaje o dropout). Hay dos excepciones declaradas.
La MLP usa una tasa inicial de $10^{-4}$, porque con $10^{-3}$ colapsó a predecir
una sola clase en el Avance~4. La variante \emph{sin regularización} retira a la
vez dropout, aumento y decaimiento de pesos (Adam en lugar de AdamW), para medir
su efecto conjunto.

\subsection{Métricas de entrenamiento comparables con validación}

Las métricas que reporta el ciclo de entrenamiento se calculan con aumento y
dropout activos, por lo que no son comparables con validación. Para medir la
brecha de generalización, en cada época se evalúan además 3,712 imágenes fijas de
entrenamiento (el mismo tamaño que validación) sin aumento y con la red en modo
inferencia. Con ellas se calculan la pérdida, la exactitud y el F1 macro de
entrenamiento. La \emph{brecha F1} es el F1 macro de ese subconjunto menos el de
validación, ambos en la época del checkpoint.

\subsection{Alcance de la evaluación}

La validación se usa para elegir el checkpoint y comparar configuraciones; sus
métricas no son una estimación independiente de generalización. Como el
conjunto de prueba de Kaggle no tiene etiquetas locales, la medida de prueba se
obtiene enviando las predicciones a la competencia
(Sección~\ref{sec:kaggle}). Cada configuración se entrenó con una sola semilla,
por lo que diferencias de pocas milésimas de F1 macro no se interpretan como
mejoras. Los resultados de este documento se generaron con TensorFlow~2.18 en una
GPU Apple M4 (\texttt{tensorflow-metal}); las corridas del Avance~4 se hicieron
con PyTorch/DirectML, por lo que sus cifras no se mezclan con las de este reporte.

\subsection{Reproducibilidad}
\label{sec:repo}

El código, las configuraciones y los resultados están en el repositorio público
\url{https://github.com/Lunaaaalj/NeuralNetworksTeamEnvironment}. El
\texttt{README} explica la instalación y los comandos. El notebook
\path{notebooks/reporte_final.ipynb} reconstruye de principio a fin la tabla de
resultados (\path{reports/final/resultados.csv}), las tablas y las figuras de
este documento, a partir de los \texttt{history.csv}, \texttt{summary.json} y
\texttt{eval\_report.json} versionados de cada experimento; con
\texttt{RUN\_TRAINING = True} también vuelve a entrenarlos. La matriz de
experimentos está en \path{scripts/run_final.py} y la configuración de referencia
en \path{configs/final.yaml}.
